% ════════════════════════════════════════════════════════════════════════════
%  Formation Engine — The Fuzzy Ball Model:
%  Multi-Layer Directional Descriptor Visualization of Coarse-Grained Beads
%  Printer-ready LaTeX.  Compile with pdflatex (twice for refs).
% ════════════════════════════════════════════════════════════════════════════
\documentclass[12pt, a4paper, oneside]{article}

% ── Geometry ────────────────────────────────────────────────────────────────
\usepackage[
  top=1.0in, bottom=1.0in, left=1.15in, right=1.15in,
  headheight=14pt
]{geometry}

% ── Fonts / encoding ───────────────────────────────────────────────────────
\usepackage[utf8]{inputenc}
\usepackage[T1]{fontenc}
\usepackage{lmodern}
\usepackage{microtype}

% ── Math ───────────────────────────────────────────────────────────────────
\usepackage{amsmath, amssymb, amsthm, mathtools}
\usepackage{bm}

% ── Tables ─────────────────────────────────────────────────────────────────
\usepackage{booktabs}
\usepackage{array}
\usepackage{tabularx}

% ── Lists ──────────────────────────────────────────────────────────────────
\usepackage{enumitem}
\setlist{nosep, leftmargin=1.5em}

% ── Code ───────────────────────────────────────────────────────────────────
\usepackage{listings}
\usepackage{xcolor}

\lstset{
  basicstyle=\ttfamily\small,
  keywordstyle=\color{blue},
  commentstyle=\color{gray},
  stringstyle=\color{red},
  breaklines=true,
  columns=fullflexible,
  keepspaces=true,
  frame=single,
  backgroundcolor=\color{gray!8},
  xleftmargin=0.5em,
  xrightmargin=0.5em,
  language=C++
}

% ── Graphics ───────────────────────────────────────────────────────────────
\usepackage{graphicx}

% ── Headers / footers ─────────────────────────────────────────────────────
\usepackage{fancyhdr}
\pagestyle{fancy}
\fancyhf{}
\fancyhead[L]{\small\textit{Formation Engine Methodology}}
\fancyhead[R]{\small\textit{Fuzzy Ball Model — Directional Descriptor Visualization}}
\fancyfoot[C]{\thepage}
\renewcommand{\headrulewidth}{0.4pt}

% ── Section formatting ─────────────────────────────────────────────────────
\usepackage{titlesec}
\titleformat{\section}
  {\Large\bfseries}
  {§\thesection}{1em}{}
\titleformat{\subsection}
  {\large\bfseries}
  {\thesubsection}{1em}{}
\titleformat{\subsubsection}
  {\normalsize\bfseries}
  {\thesubsubsection}{1em}{}

% ── Hyperlinks ─────────────────────────────────────────────────────────────
\usepackage[
  colorlinks=true, linkcolor=black, citecolor=black, urlcolor=blue
]{hyperref}

% ── Convenience macros ─────────────────────────────────────────────────────
\newcommand{\file}[1]{\texttt{#1}}
\newcommand{\type}[1]{\texttt{#1}}
\newcommand{\field}[1]{\texttt{#1}}
\newcommand{\Ang}{\text{\AA}}
\newcommand{\kcalmol}{\text{kcal/mol}}
\newcommand{\SH}{Y_\ell^m}
\newcommand{\Rvec}{\bm{r}}
\newcommand{\nhat}{\hat{\bm{n}}}

% ════════════════════════════════════════════════════════════════════════════
%  Title
% ════════════════════════════════════════════════════════════════════════════
\title{%
  \textbf{Formation Engine Methodology}\\[6pt]
  \textbf{The Fuzzy Ball Model}\\[10pt]
  \large Multi-Layer Directional Descriptor Visualization\\
  of Coarse-Grained Atomistic Beads\\[4pt]
  \normalsize Spherical Harmonic Surface Fields, Anisotropic Shell Rendering,
              and Environment-Responsive Visual Inspection
}
\author{%
  % ── Replace with your name, institution, and email ──
  VSEPR-SIM Project%
  % \\ \small Your Institution
  % \\ \small \texttt{your.email@example.com}
}
\date{Version 3.0.0 — \today}

% ════════════════════════════════════════════════════════════════════════════
\begin{document}
\maketitle

% ════════════════════════════════════════════════════════════════════════════
%  Abstract
% ════════════════════════════════════════════════════════════════════════════
\begin{abstract}
This document describes the \emph{Fuzzy Ball Model}, a multi-layer visual
representation framework for coarse-grained (CG) beads in the Formation
Engine.  Rather than treating each CG bead as a featureless point mass,
the fuzzy ball model preserves and exposes the directional, anisotropic,
and environment-responsive character of the underlying atomistic fragment.
Each bead is rendered as a translucent spherical shell whose effective
radius encodes atomistic extent, whose surface carries a three-channel
descriptor field expanded in real spherical harmonics, and whose interior
optionally reveals the source atomistic skeleton.  Eight composable view
modes---\emph{Shell}, \emph{Scaffold}, \emph{SurfaceState},
\emph{InternalRef}, \emph{Skeleton}, \emph{Cutaway}, \emph{Residual}, and
\emph{Comparison}---permit progressive disclosure from macroscopic overview
to fragment-level inspection.  A five-level scale hierarchy
(\emph{SceneCloud}~$\to$~\emph{FragmentCompare}) governs level-of-detail
transitions.  The model is deterministic, fully inspectable, and exports
structured data suitable for scientific reporting.  Implementation:
\file{coarse\_grain/vis/bead\_visual\_record.hpp} and
\file{coarse\_grain/vis/bead\_inspector\_view.hpp}.
\end{abstract}

\tableofcontents
\newpage

% ════════════════════════════════════════════════════════════════════════════
\section{Introduction}
% ════════════════════════════════════════════════════════════════════════════

Coarse-grained molecular simulation replaces groups of atoms with
single interaction sites (beads), reducing degrees of freedom by one
to two orders of magnitude.  The standard visual representation of a
CG bead is a uniform sphere, which discards all information about
internal geometry, anisotropy, and interaction directionality.

The \textbf{Fuzzy Ball Model} restores this information by treating
each bead not as a hard marble but as a \emph{semi-transparent shell
with directional structure}---a ``fuzzy ball.''  The model achieves
this through three complementary data layers:

\begin{enumerate}
  \item \textbf{Layer 1 --- Effective Shell.}  A spherical shell whose
        radius, opacity, colour, and deformation encode the bead's
        effective extent, mapping confidence, chemical class, and
        principal anisotropy.
  \item \textbf{Layer 2 --- Surface-State Field.}  A structured
        $(\theta, \varphi)$ grid of steric, electrostatic, and
        dispersion descriptor values painted as a spherical heatmap
        over the shell.
  \item \textbf{Layer 3 --- Internal Structural Reference.}  The
        original atomistic fragment (atoms, bonds, orientation markers)
        rendered inside the translucent shell for direct comparison.
\end{enumerate}

These layers are combined through eight selectable \emph{view modes}
that control which layers are active and at what opacity.  A five-level
\emph{scale hierarchy} determines how much detail is rendered at each
zoom level, from a distant scene cloud to a close-up fragment comparison.

The model is implemented entirely within the
\file{coarse\_grain/vis/} subsystem and consumes the
\type{BeadVisualRecord} adapter object produced by
\texttt{build\_visual\_record()}.

% ════════════════════════════════════════════════════════════════════════════
\section{Theoretical Foundation}
% ════════════════════════════════════════════════════════════════════════════

\subsection{Effective Radius and Anisotropy}

Each CG bead maps a molecular fragment to a single centre $\Rvec_B$.
The \emph{effective radius} $R_{\text{eff}}$ quantifies the spatial
extent of this fragment.  When a surface descriptor probe radius is
available,
\begin{equation}
  R_{\text{eff}} = r_{\text{probe}}, \qquad r_{\text{probe}} > 0,
  \label{eq:reff_probe}
\end{equation}
otherwise a mass-based fallback is used:
\begin{equation}
  R_{\text{eff}} = 1.5 \left(\frac{M_B}{12}\right)^{1/3}, \qquad
  R_{\text{eff}} \in [0.5, \, 4.0] \;\Ang.
  \label{eq:reff_mass}
\end{equation}

The bead is not spherically symmetric in general.  The inertia tensor
of the source fragment yields three principal eigenvalues
$\lambda_1 \leq \lambda_2 \leq \lambda_3$.  Their ratios define an
\emph{anisotropy vector},
\begin{equation}
  \bm{a}_B = (\lambda_1, \, \lambda_2, \, \lambda_3),
  \label{eq:aniso}
\end{equation}
which deforms the shell rendering from a perfect sphere toward an
ellipsoid.

\subsection{Local Frame Assignment}

A local coordinate frame $\{\hat{\bm{e}}_1, \hat{\bm{e}}_2,
\hat{\bm{e}}_3\}$ is assigned to each bead via one of several
strategies, selected automatically based on fragment topology:

\begin{table}[htbp]
  \centering
  \caption{Frame-method selection rules.}
  \label{tab:frame_methods}
  \begin{tabular}{lll}
    \toprule
    Method & Condition & Description \\
    \midrule
    \texttt{InertiaFit}        & Full surface data available
      & PCA of inertia tensor \\
    \texttt{RingPlaneFit}      & Fragment is cyclic / aromatic
      & $\hat{\bm{e}}_3$ = ring-plane normal \\
    \texttt{MetalCoordination} & Fragment is organometallic
      & Principal axis of coordination shell \\
    \texttt{PCA\_Fallback}     & Orientation normal only
      & Gram--Schmidt from $\hat{\bm{n}}$ \\
    \texttt{SingleAtom}        & Fragment has $\leq 1$ atom
      & Degenerate; confidence $= 0$ \\
    \bottomrule
  \end{tabular}
\end{table}

Frame confidence $\kappa_B$ is computed from the eigenvalue
separation:
\begin{equation}
  \kappa_B = \min\!\left(1, \;
    3 \, \frac{\lambda_3 - \lambda_1}
              {\lambda_1 + \lambda_2 + \lambda_3}\right).
  \label{eq:frame_confidence}
\end{equation}
The rendering line-width of the frame triad axes is set to 3.0 when
$\kappa_B > 0.5$ and 1.5 otherwise, providing a visual cue for
confidence.

\subsection{Mapping Confidence and Residual}

The mapping from atomistic fragment to CG bead introduces a
positional error $\delta = |\Rvec_{\text{COM}} - \Rvec_{\text{COG}}|$.
The mapping confidence is
\begin{equation}
  c_B = \max\!\left(0, \; 1 - \frac{\delta}{2}\right),
  \label{eq:mapping_confidence}
\end{equation}
which scales shell opacity: $\alpha_{\text{shell}} =
\alpha_0 \, (0.3 + 0.7 \, c_B)$.  Poorly mapped beads become
visually faint, drawing the researcher's attention to reliable regions.

% ════════════════════════════════════════════════════════════════════════════
\subsection{Correspondence with the Canonical Container Geometry}
\label{sec:container-correspondence}
% ════════════════════════════════════════════════════════════════════════════

The Identity--State Decomposition Framework (§0) defines for every
particle a \emph{Canonical Container Geometry Approximation}:
%
\begin{equation}
  \Omega_i =
  \bigl(
    R_i,\;
    \mathcal{T}_i(\theta,\phi),\;
    \rho_i(r)
  \bigr)
  \tag{§0, Eq.\,\ref*{eq:envelope}$'$}
  \label{eq:container-omega}
\end{equation}
%
Each component of~$\Omega_i$ has a direct, implemented correspondence
in the Fuzzy Ball Model.

\paragraph{Effective radius $R_i \;\longleftrightarrow\;
            R_{\mathrm{eff}}$.}
The §0 definition,
\(R_i = \alpha_Z\,Z_i^{1/3} + \alpha_Q\,|Q_i|\),
provides a \emph{first-principles} estimate suitable for isolated
atoms.  In the CG context, the bead aggregates many atoms, so the
Fuzzy Ball Model uses a \emph{probe-radius} or mass-derived fallback
(Equations~\ref{eq:reff_probe}--\ref{eq:reff_mass}).  Both converge
to the same role: determining the shell extent displayed by Layer~1.
The §0 radius governs the formal interaction envelope; the
visualisation radius governs what the researcher sees.  They share
the invariant:
\begin{equation}
  R_{\text{eff}} \;\propto\; R_i \;\propto\; V_{\text{frag}}^{1/3},
  \label{eq:radius-link}
\end{equation}
where $V_{\text{frag}}$ is the volumetric extent of the mapped
fragment.

\paragraph{Angular texture $\mathcal{T}_i(\theta,\phi)
  \;\longleftrightarrow\; S^{(k)}(\theta,\varphi)$.}
The §0 angular texture is a single spherical-harmonic expansion that
encodes the directional interaction shape of the canonical container.
The Fuzzy Ball Model implements a \emph{three-channel} generalisation:
%
\begin{equation}
  \mathcal{T}_i(\theta,\phi) \;\longrightarrow\;
  \bigl\{
    S^{(0)},\; S^{(1)},\; S^{(2)}
  \bigr\}(\theta,\varphi),
  \label{eq:texture-link}
\end{equation}
%
one expansion per physical channel (steric, electrostatic, dispersion).
The synthesised channel $S^{(\text{syn})}$ (Equation~\ref{eq:synthesized})
recovers the single-scalar §0 texture as a weighted contraction.
The implementation uses the same mathematical object---real spherical
harmonics $Y_\ell^m$---at both levels, so the correspondence is
exact up to the number of channels.

\paragraph{Radial density envelope
  $\rho_i(r) \;\longleftrightarrow\;$ radial exaggeration $\xi$.}
The §0 radial envelope,
\(\rho_i(r) = \rho_0\,e^{-r/\lambda_i}\),
describes a soft exponential decay of the interaction field.
In the Fuzzy Ball Model, this manifests as the \emph{radial
exaggeration} parameter~$\xi$
(Equation~\ref{eq:radial_exag}), which deforms the rendered surface
proportionally to the local field strength.  When $\xi > 0$ the
shell bulges outward in directions of high $S^{(k)}$, reproducing
the intuition that $\rho_i(r)$ is larger in directions where the
interaction field is strong.  The formal mapping is:
\begin{equation}
  \rho_i\!\bigl(R_{\text{eff}} \cdot \nhat(\theta,\varphi)\bigr)
  \;\sim\;
  1 + \xi \, \bar{S}^{(k)}(\theta,\varphi),
  \label{eq:density-link}
\end{equation}
providing a visual analogue of the radial density profile.

\paragraph{Summary.}
The correspondence is collected in Table~\ref{tab:omega-mapping}.

\begin{table}[htbp]
  \centering
  \caption{Mapping between §0 Canonical Container Geometry
           $\Omega_i$ and the Fuzzy Ball Model implementation.}
  \label{tab:omega-mapping}
  \begin{tabular}{lll}
    \toprule
    §0 Component & Fuzzy Ball Implementation & Source \\
    \midrule
    $R_i$ (effective radius)
      & \field{effective\_radius}
      & \file{bead\_visual\_record.hpp} \\
    $\mathcal{T}_i(\theta,\phi)$ (angular texture)
      & $S^{(k)}(\theta,\varphi)$ via SH expansion
      & \type{SurfaceStateGrid} \\
    $\rho_i(r)$ (radial density)
      & Radial exaggeration $\xi$
      & \field{radial\_exaggeration} \\
    $\Sigma_i$ (structural role)
      & \field{structural\_role}
      & \type{StructuralRole} enum \\
    $\Lambda_i$ (stability class)
      & \field{stability\_class}
      & \type{StabilityClass} enum \\
    \bottomrule
  \end{tabular}
\end{table}

% ════════════════════════════════════════════════════════════════════════════
\subsection{Structural Role and Stability Class}
\label{sec:role-stability}
% ════════════════════════════════════════════════════════════════════════════

The Identity--State Decomposition Framework assigns each particle
two discrete classification fields: a \emph{structural role}
$\Sigma_i$ and a \emph{stability class}~$\Lambda_i$.  Both are now
carried on every \type{Bead} and propagated into the
\type{BeadVisualRecord} for rendering-time inspection.

\paragraph{Structural role $\Sigma_i$.}
A five-valued prior encoding dominant bonding topology:
%
\[
  \Sigma_i \in
  \bigl\{\,
    0\!=\!\text{inert},\;
    1\!=\!\text{ionic},\;
    2\!=\!\text{directional covalent},\;
    3\!=\!\text{metallic},\;
    4\!=\!\text{mixed}
  \,\bigr\}.
\]
%
The role biases the three-channel interaction kernel via per-channel
weights $w_k^{(\Sigma)}$ (Table~\ref{tab:role-weights}).  For a
bead pair $(A,B)$ the combined weight uses a geometric mean,
analogous to Lorentz--Berthelot mixing:
%
\begin{equation}
  w_k^{AB} = \sqrt{w_k^A \cdot w_k^B}.
  \label{eq:role-mix}
\end{equation}

\begin{table}[htbp]
  \centering
  \caption{Structural role $\Sigma_i$: per-channel weight biases.
           Values are multipliers on the existing coupling constants
           $\lambda_k$.}
  \label{tab:role-weights}
  \begin{tabular}{lccc}
    \toprule
    $\Sigma$ & $w_{\text{steric}}$ & $w_{\text{elec}}$ & $w_{\text{disp}}$ \\
    \midrule
    0 — Inert              & 0.3 & 0.1 & 1.5 \\
    1 — Ionic-Dominant     & 0.5 & 1.5 & 0.3 \\
    2 — Directional Cov.   & 1.5 & 0.5 & 0.7 \\
    3 — Metallic           & 1.0 & 0.2 & 1.3 \\
    4 — Mixed              & 1.0 & 1.0 & 1.0 \\
    \bottomrule
  \end{tabular}
\end{table}

\paragraph{Stability class $\Lambda_i$.}
A four-valued persistence classification:
\[
  \Lambda_i \in
  \bigl\{\,
    0\!=\!\text{transient},\;
    1\!=\!\text{metastable},\;
    2\!=\!\text{ambient-stable},\;
    3\!=\!\text{bulk-lattice}
  \,\bigr\}.
\]
Both fields are visible in the Comparison view mode and in the ImGui
inspector panel, satisfying the anti-black-box requirement that every
identity-state component of a bead be directly inspectable.

Implementation: \type{StructuralRole} and \type{StabilityClass} enums
in \file{coarse\_grain/core/bead.hpp}; the classifier
\texttt{classify\_structural\_role()} infers $\Sigma_i$ from the
$Z$~values of the parent atom group.

% ════════════════════════════════════════════════════════════════════════════
\section{Layer 2: Surface-State Descriptor Field}
\label{sec:layer2}
% ════════════════════════════════════════════════════════════════════════════

\subsection{Spherical Harmonic Expansion}

Each bead carries a \emph{unified descriptor} consisting of three
scalar fields on $S^2$, one per interaction channel $k$:
\begin{equation}
  S^{(k)}(\theta, \varphi)
    = \sum_{\ell=0}^{\ell_{\max}}
      \sum_{m=-\ell}^{\ell}
      c_{\ell m}^{(k)} \, \SH(\theta, \varphi),
  \label{eq:sh_expansion}
\end{equation}
where $k \in \{0 = \text{steric},\; 1 = \text{electrostatic},\;
2 = \text{dispersion}\}$ and $\SH$ are real spherical harmonics.

The coefficient vectors $\{c_{\ell m}^{(k)}\}$ are pre-solved during
bead construction and stored in the \type{UnifiedDescriptor}.  At
visualisation time they are evaluated on a structured angular grid.

\subsection{Surface-State Grid}

The Layer~2 grid tessellates $S^2$ into
$(N_\theta + 1) \times N_\phi$ vertices in row-major order.  The
default resolution is $N_\theta = 24$, $N_\phi = 48$, yielding
1\,200 vertices:

\begin{equation}
  \theta_i = \frac{\pi \, i}{N_\theta}, \quad i = 0 \ldots N_\theta,
  \qquad
  \varphi_j = \frac{2\pi \, j}{N_\phi}, \quad j = 0 \ldots N_\phi - 1.
  \label{eq:grid_coords}
\end{equation}

At each vertex, all three channels are evaluated via
Equation~\ref{eq:sh_expansion}.  A fourth \emph{synthesised} channel
combines them:
\begin{equation}
  S^{(\text{syn})}(\theta, \varphi)
    = w_s \, S^{(0)} + w_e \, S^{(1)} + w_d \, S^{(2)},
  \label{eq:synthesized}
\end{equation}
with equal default weights $w_s = w_e = w_d = \tfrac{1}{3}$.

Per-channel value ranges $[S_{\min}^{(k)}, S_{\max}^{(k)}]$ are
tracked for colourmap normalisation.  The active channel is selectable
at runtime.

\subsection{Surface-State Channels}

\begin{table}[htbp]
  \centering
  \caption{Surface-state descriptor channels.}
  \label{tab:channels}
  \begin{tabular}{clp{7.5cm}}
    \toprule
    $k$ & Name & Physical Meaning \\
    \midrule
    0 & Steric        & Geometric accessibility / excluded-volume envelope \\
    1 & Electrostatic & Charge-field distribution from partial charges \\
    2 & Dispersion    & Attractive interaction strength (London dispersion) \\
    syn & Synthesized  & Weighted combination of all three channels \\
    \bottomrule
  \end{tabular}
\end{table}

\subsection{Rendering the Heatmap}

The grid is rendered as a quad-strip tessellation over the shell
surface.  Each vertex colour is obtained by mapping the active
channel's value through a sequential colourmap, normalised by the
tracked min/max range.

Optional \emph{radial exaggeration} $\xi$ deforms the sphere in the
normal direction:
\begin{equation}
  \Rvec(\theta, \varphi)
    = \Rvec_B + R_{\text{eff}} \,
      \bigl(1 + \xi \, \bar{S}^{(k)}(\theta,\varphi)\bigr) \,
      \nhat(\theta, \varphi),
  \label{eq:radial_exag}
\end{equation}
where $\bar{S}^{(k)}$ is the normalised channel value in $[0,1]$.
When $\xi = 0$ the surface remains a perfect sphere; larger values
produce a ``bumpy'' surface proportional to the field strength.

% ════════════════════════════════════════════════════════════════════════════
\section{Layer 3: Internal Structural Reference}
% ════════════════════════════════════════════════════════════════════════════

Layer~3 renders the atomistic fragment that was mapped into the bead.
This allows direct visual comparison between the CG representation and
its atomistic source.

\subsection{Source Fragment Miniature}

The \type{SourceFragmentMiniature} stores a lightweight copy of the
atomistic geometry:

\begin{itemize}
  \item \textbf{Atoms}: position, atomic number, Van der Waals radius
        (simplified lookup), and topological flags.
  \item \textbf{Bonds}: atom index pair $(i,j)$ and bond order.
  \item \textbf{Centres}: centre-of-mass $\Rvec_{\text{COM}}$ and
        centre-of-geometry $\Rvec_{\text{COG}}$.
\end{itemize}

\subsection{Internal Reference Rendering Options}

\begin{table}[htbp]
  \centering
  \caption{Configurable options for Layer~3 rendering.}
  \label{tab:internal_ref}
  \begin{tabular}{llr}
    \toprule
    Option & Description & Default \\
    \midrule
    \field{show\_atoms}              & Draw atom spheres         & true  \\
    \field{show\_bonds}              & Draw bond sticks          & true  \\
    \field{show\_orientation\_markers}& Draw local frame axes    & true  \\
    \field{element\_coloring}        & CPK element colours       & true  \\
    \field{atom\_scale}              & Fraction of VdW radius    & 0.25  \\
    \field{bond\_width}              & Bond line width (px)      & 1.5   \\
    \bottomrule
  \end{tabular}
\end{table}

% ════════════════════════════════════════════════════════════════════════════
\section{View Modes}
\label{sec:viewmodes}
% ════════════════════════════════════════════════════════════════════════════

The eight view modes control the composition of Layers~1--3 and
additional diagnostic overlays.

\begin{table}[htbp]
  \centering
  \caption{View-mode composition matrix.  $\alpha$ denotes shell
           opacity.  ``Triad'' = local frame triad.  ``Heatmap'' =
           Layer~2 surface-state field.  ``Fragment'' = Layer~3
           atomistic reference.  ``Halo'' = residual glyph.}
  \label{tab:viewmodes}
  \begin{tabular}{lcccccc}
    \toprule
    Mode & $\alpha$ & Triad & Scaffold & Heatmap & Fragment & Halo \\
    \midrule
    Shell        & 0.60 &     &     &     &     &     \\
    Scaffold     & 0.25 & \checkmark & \checkmark &     &     &     \\
    SurfaceState & 0.10 & \checkmark &     & \checkmark &     &     \\
    InternalRef  & 0.12 & \checkmark &     &     & \checkmark &     \\
    Skeleton     & 0.15 & \checkmark &     &     & \checkmark &     \\
    Cutaway      & 0.50 & \checkmark & \checkmark &     &     &     \\
    Residual     & 0.30 & \checkmark &     &     &     & \checkmark \\
    Comparison   & 0.15 & \checkmark &     & \checkmark & \checkmark & \checkmark \\
    \bottomrule
  \end{tabular}
\end{table}

\paragraph{Shell.}
The simplest mode.  A semi-transparent sphere at
$R_{\text{eff}}$, colour-coded by chemical class:
purple for aromatics, gold for metals, teal for cyclic fragments,
blue for all others.  Opacity encodes mapping confidence.

\paragraph{Scaffold.}
Adds the solved direction vectors (coordination geometry) as
radial lines from the bead centre, plus the surface descriptor
patch field rendered as directional glyphs.

\paragraph{SurfaceState.}
The Layer~2 heatmap mode.  The shell becomes a ghost ($\alpha = 0.10$)
and is overlaid with the spherical descriptor field.  The active channel
is selectable (steric, electrostatic, dispersion, or synthesised).

\paragraph{InternalRef.}
The Layer~3 mode.  The atomistic fragment is rendered inside the
ghost shell using ball-and-stick representation with CPK colouring.

\paragraph{Skeleton.}
Similar to InternalRef but displays the raw stick model of the source
fragment.

\paragraph{Cutaway.}
A clipping plane slices the shell to expose internal structure.  The
scaffold and descriptor patches are drawn within the visible
hemisphere.

\paragraph{Residual.}
Displays a residual halo whose size and colour encode the descriptor
reconstruction residual, geometry distortion, and frame ambiguity.
Used to identify beads where the CG mapping is least faithful.

\paragraph{Comparison.}
All diagnostic layers simultaneously: ghost shell, heatmap, atomistic
fragment, frame triad, and residual halo.  The most information-dense
mode, intended for close-up single-bead analysis.

% ════════════════════════════════════════════════════════════════════════════
\section{Scale Hierarchy}
% ════════════════════════════════════════════════════════════════════════════

Five scale levels govern the level of detail rendered for the bead
population as a whole.

\begin{table}[htbp]
  \centering
  \caption{Scale hierarchy levels.}
  \label{tab:scale}
  \begin{tabular}{cll}
    \toprule
    Level & Name & Description \\
    \midrule
    0 & SceneCloud       & Distant overview; beads as point sprites \\
    1 & Cluster          & Intermediate; smooth shaded spheres \\
    2 & BeadDetail       & Full shell rendering with opacity/colour \\
    3 & BeadInterior     & Internal structure visible through ghost shell \\
    4 & FragmentCompare  & Single-bead comparison with atomistic overlay \\
    \bottomrule
  \end{tabular}
\end{table}

Transitions between levels are governed by camera distance and may be
overridden for selected beads via the inspector panel.

% ════════════════════════════════════════════════════════════════════════════
\section{Alpha Polarisability Model Integration}
\label{sec:alpha}
% ════════════════════════════════════════════════════════════════════════════

The ``fuzzy'' character of the ball extends beyond visualisation into
the physical model.  The \emph{alpha model} fits atomic polarisabilities
using a parameterised form that includes \emph{blob} terms:

\begin{equation}
  \alpha_{\text{eff}} = f_{\text{blob}}(\alpha_0, \, Z, \, \text{block}),
  \label{eq:alpha_blob}
\end{equation}
where $f_{\text{blob}}$ is controlled by three fitted parameters:
\begin{itemize}
  \item \field{blob\_f\_lin} --- linear fuzzy-shell mixing fraction,
  \item \field{blob\_f\_half} --- half-shell contribution,
  \item \field{blob\_f\_full} --- full-shell contribution.
\end{itemize}

These are determined by evolutionary stochastic optimisation against
reference polarisability data (see \file{tools/fit\_alpha\_model.cpp}).
Additional fitted parameters include:

\begin{table}[htbp]
  \centering
  \caption{Alpha model parameters (Method~D).}
  \label{tab:alpha_params}
  \begin{tabular}{lp{8.5cm}}
    \toprule
    Parameter & Description \\
    \midrule
    $k_{\text{period}}[1\ldots7]$ & Period-dependent scaling factors \\
    $c_{\text{block}}[s,p,d,f]$   & Orbital-block correction terms \\
    $c_{\text{rel}}$              & Relativistic correction coefficient \\
    $\beta_f$                     & $f$-block contraction exponent \\
    $a_{f1}, a_{f2}, a_{f3}$      & Lanthanide/actinide amplitude terms \\
    $\sigma_{f1}, \sigma_{f2}$    & Gaussian widths for $f$-block \\
    $b_{\text{bind}}$             & Binding-energy correction \\
    $q_{\text{drude}}$            & Drude oscillator effective charge \\
    \bottomrule
  \end{tabular}
\end{table}

The blob parameters connect the visual ``fuzziness'' of the shell to
a quantitative physical property: the polarisability of the electron
cloud determines how ``soft'' and ``diffuse'' each bead appears, both
in the force field and in the rendering.

% ════════════════════════════════════════════════════════════════════════════
\section{Environment-Responsive Coupling}
% ════════════════════════════════════════════════════════════════════════════

The fuzzy ball adapts to its thermodynamic environment.  The
\type{BeadVisualRecord} captures an environment-state snapshot at each
frame:

\begin{table}[htbp]
  \centering
  \caption{Environment observables stored per bead.}
  \label{tab:env}
  \begin{tabular}{clp{7cm}}
    \toprule
    Symbol & Field & Meaning \\
    \midrule
    $\eta_B$ & \field{eta} & Slow internal state (relaxation variable) \\
    $\rho_B$ & \field{rho} & Local bead density \\
    $f_B$    & \field{target\_f} & Target function value \\
    $P_{2,B}$& \field{P2}  & Nematic order parameter \\
    $C_B$    & \field{coordination\_C} & Coordination number \\
    \bottomrule
  \end{tabular}
\end{table}

These observables modulate the three-channel kernel described in
the Seed \& Bead specification
($\gamma_{\text{steric}}$, $\gamma_{\text{elec}}$,
$\gamma_{\text{disp}}$), closing the loop between simulation dynamics
and visual representation.

% ════════════════════════════════════════════════════════════════════════════
\section{Connectivity and Neighbour Interactions}
% ════════════════════════════════════════════════════════════════════════════

Each bead record stores a list of \type{NeighbourLink} entries that
describe pairwise interactions with adjacent beads:

\begin{itemize}
  \item \field{neighbour\_bead\_id} --- index of the neighbour bead,
  \item \field{distance} --- centre-to-centre separation ($\Ang$),
  \item \field{interaction\_strength} --- magnitude of the pairwise
        potential,
  \item \field{interaction\_type} --- string label classifying the
        geometry (e.g.\ ``face-to-face'', ``edge-to-face'').
\end{itemize}

This information enables graph-based analyses and can be exported to
the snapshot reporting system for automated figure generation.

% ════════════════════════════════════════════════════════════════════════════
\section{Implementation Summary}
% ════════════════════════════════════════════════════════════════════════════

\subsection{Data Flow}

The conversion from simulation data to visual output follows a
deterministic pipeline:

\begin{enumerate}
  \item \textbf{Bead construction.}  The atomistic fragment is mapped
        to a CG bead via the \texttt{6+9 stepper}, producing a
        \type{Bead} with position, mass, surface descriptor, and
        unified descriptor.
  \item \textbf{Record building.}  \texttt{build\_visual\_record()}
        converts the \type{Bead} into a \type{BeadVisualRecord},
        computing effective radius, local frame, anisotropy,
        descriptor grids, and source-fragment miniature.
  \item \textbf{Inspector rendering.}  \type{BeadInspectorView}
        consumes the record and dispatches to the selected view mode,
        composing Layers~1--3 with the appropriate opacity and overlay
        settings.
\end{enumerate}

\subsection{Key Source Files}

\begin{table}[htbp]
  \centering
  \caption{Implementation files for the fuzzy ball model.}
  \label{tab:files}
  \begin{tabular}{lp{8cm}}
    \toprule
    File & Role \\
    \midrule
    \file{coarse\_grain/vis/bead\_visual\_record.hpp}
      & Data model: \type{BeadVisualRecord}, enumerations,
        grid types, \texttt{build\_visual\_record()} \\
    \file{coarse\_grain/vis/bead\_inspector\_view.hpp}
      & Renderer: view-mode dispatch, shell/scaffold/heatmap/
        fragment/residual compositing \\
    \file{tools/fit\_alpha\_model.cpp}
      & Alpha polarisability fitting with blob parameters \\
    \file{tools/auto\_fit\_alpha.cpp}
      & Evolutionary optimiser for \type{AlphaModelParams} \\
    \file{coarse\_grain/models/seed\_bead\_stepper.hpp}
      & 6+9 step function that produces bead data consumed by
        the visual record builder \\
    \bottomrule
  \end{tabular}
\end{table}

\subsection{Descriptor Evaluation: Spherical Harmonics}

The inline evaluator computes real spherical harmonics $\SH(\theta,\varphi)$
up to $\ell_{\max}$ and contracts with the coefficient vector:

\begin{lstlisting}
inline double evaluate_sh_expansion_at(
    const std::vector<double>& coeffs,
    int l_max,
    double theta, double phi)
{
    auto Y = evaluate_all_harmonics_dynamic(theta, phi, l_max);
    double val = 0.0;
    int n = std::min((int)coeffs.size(), (int)Y.size());
    for (int i = 0; i < n; ++i)
        val += coeffs[i] * Y[i];
    return val;
}
\end{lstlisting}

% ════════════════════════════════════════════════════════════════════════════
\section{Results and Discussion}
% ════════════════════════════════════════════════════════════════════════════

The fuzzy ball model provides several concrete improvements over
conventional CG visualisation:

\begin{enumerate}
  \item \textbf{Directional information is preserved.}  The three-channel
        surface-state heatmap reveals steric, electrostatic, and dispersion
        anisotropy that would be invisible with uniform spheres.

  \item \textbf{Mapping quality is visually encoded.}  Shell opacity
        and residual halos immediately flag beads whose CG representation
        deviates significantly from the atomistic source.

  \item \textbf{Multi-scale inspection is seamless.}  The five-level
        scale hierarchy and eight view modes allow a researcher to
        navigate from a macroscopic scene overview to single-fragment
        comparison without switching tools.

  \item \textbf{Physical fuzziness is quantified.}  The alpha-model
        blob parameters tie the visual softness of each shell to a
        fitted atomic polarisability, connecting rendering to a
        measurable physical property.

  \item \textbf{Environment responsiveness is captured.}  The
        environment-state snapshot ($\eta_B$, $\rho_B$, $P_{2,B}$,
        $C_B$, $f_B$) modulates the descriptor kernel, ensuring the
        visual representation reflects local thermodynamic conditions.
\end{enumerate}

% ════════════════════════════════════════════════════════════════════════════
\section{Conclusion}
% ════════════════════════════════════════════════════════════════════════════

The Fuzzy Ball Model transforms the conventional point-sphere
representation of coarse-grained beads into a scientifically rich,
multi-layer visual instrument.  By encoding directional descriptor
fields via spherical harmonic expansion, exposing the internal
atomistic skeleton through progressive transparency, and coupling
shell appearance to environment state and fitted polarisability
parameters, the model enables deterministic, inspectable, and
research-grade visualisation.

The eight-mode composition system and five-level scale hierarchy
provide a systematic framework for navigating between overview and
detail, supporting both qualitative exploration and quantitative
analysis.  All outputs are fully deterministic and exportable for
inclusion in scientific reports.

Future work includes adaptive tessellation based on descriptor
gradient magnitude, time-dependent surface-state animation for
trajectory playback, and integration of the fuzzy ball rendering
with the automated snapshot and graph reporting pipeline.

% ════════════════════════════════════════════════════════════════════════════
%  References
% ════════════════════════════════════════════════════════════════════════════
\bibliographystyle{plain}
% \bibliography{references}  % Uncomment if you have a .bib file

\end{document}
